\subsubsection{Computational Complexity and Hardness}

Computational complexity theory formalizes the inherent difficulty of solving problems and provides a framework for classifying problems by resource requirements. The primary resource of interest is time complexity—the number of computational steps required to solve a problem as a function of input size.

\paragraph{Complexity Classes and Polynomial-Time Reductions}

Problems are classified into complexity classes based on the computational resources required for their solution. The class $\mathcal{P}$ (Polynomial-time) contains decision problems solvable by a deterministic Turing machine in polynomial time. The class $\mathcal{NP}$ (Nondeterministic Polynomial-time) contains decision problems whose solutions can be verified in polynomial time. Formally, a decision problem belongs to $\mathcal{NP}$ if a proposed solution can be checked for correctness in time polynomial in the input size.

A fundamental relationship exists between $\mathcal{P}$ and $\mathcal{NP}$: if $\mathcal{P} = \mathcal{NP}$, then every problem whose solution can be verified in polynomial time can also be solved in polynomial time. The $\mathcal{P}$ vs. $\mathcal{NP}$ question remains unsolved, representing a central open problem in theoretical computer science.

A problem $\Pi_1$ is polynomial-time reducible to problem $\Pi_2$ (denoted $\Pi_1 \leq_p \Pi_2$) if there exists a polynomial-time algorithm transforming any instance of $\Pi_1$ into an equivalent instance of $\Pi_2$. This reduction relationship enables problem classification: if $\Pi_1 \leq_p \Pi_2$ and $\Pi_1$ is hard, then $\Pi_2$ is at least as hard.

\paragraph{NP-Completeness and NP-Hardness}

A problem is $\mathcal{NP}$-complete if it belongs to $\mathcal{NP}$ and every problem in $\mathcal{NP}$ can be polynomial-time reduced to it. Cook's theorem established that the Boolean satisfiability problem (SAT) is $\mathcal{NP}$-complete. Subsequently, thousands of problems have been proven $\mathcal{NP}$-complete through polynomial-time reductions from known $\mathcal{NP}$-complete problems.

A problem is $\mathcal{NP}$-hard if every problem in $\mathcal{NP}$ can be polynomial-time reduced to it. $\mathcal{NP}$-hardness is a stronger condition than $\mathcal{NP}$-completeness: $\mathcal{NP}$-hard problems need not be in $\mathcal{NP}$ themselves. Optimization problems, which require finding the best solution rather than verifying a solution's validity, are typically $\mathcal{NP}$-hard.

The presumption that $\mathcal{P} \neq \mathcal{NP}$ implies that $\mathcal{NP}$-hard problems cannot be solved by polynomial-time algorithms. This fundamental barrier motivates algorithmic approaches beyond exact polynomial-time computation.

\paragraph{Implications for Vehicle Routing Problems}

The Traveling Salesman Problem (TSP), asking for the minimum-cost tour visiting all cities exactly once, is $\mathcal{NP}$-hard. Extensions introducing multiple vehicles (Vehicle Routing Problem), time windows, and heterogeneous fleet constraints remain $\mathcal{NP}$-hard. The Coordinated Truck-Drone Routing Problem belongs to this complexity class, as it generalizes the $\mathcal{NP}$-hard VRP with additional coordination constraints.

Formally, if the decision version of routing problem $\Pi_r$ (``Is there a solution with cost $\leq C$?'') belongs to $\mathcal{NP}$, and TSP is polynomial-time reducible to $\Pi_r$, then $\Pi_r$ is $\mathcal{NP}$-complete. The optimization version asking for the minimum-cost solution is $\mathcal{NP}$-hard.

\paragraph{Approximation Algorithms and Approximation Ratios}

Given the intractability of $\mathcal{NP}$-hard problems, approximation algorithms provide polynomial-time solutions with guaranteed performance bounds. An approximation algorithm for a minimization problem achieves approximation ratio $\rho$ if:

\begin{equation}
\frac{\text{ALG}(I)}{\text{OPT}(I)} \leq \rho
\end{equation}

for all instances $I$, where ALG$(I)$ is the algorithm's solution cost and OPT$(I)$ is the optimal solution cost. A $\rho$-approximation algorithm guarantees that its solution is at most $\rho$ times worse than optimal.

For the metric TSP (where distances satisfy the triangle inequality), a 1.5-approximation algorithm exists via minimum spanning tree construction. However, many routing variants lack known polynomial-time approximation schemes with good ratios, necessitating heuristic and metaheuristic approaches.

\paragraph{Hardness of Approximation}

Beyond showing that exact polynomial-time solutions are unlikely, complexity theory also reveals limits on approximation quality achievable in polynomial time. Some $\mathcal{NP}$-hard problems admit hardness of approximation results: proving that polynomial-time algorithms achieving certain approximation ratios would imply $\mathcal{P} = \mathcal{NP}$. These results establish fundamental barriers on what polynomial-time algorithms can guarantee.

The existence of $\mathcal{NP}$-hardness and approximation barriers justifies the use of heuristic, metaheuristic, and learning-based approaches for routing problems. These methods trade worst-case polynomial-time guarantees for practical solution quality, enabling effective real-world problem solving despite theoretical intractability.
